Userspace runtimes that broker untrusted code on commodity Linux are
stuck choosing between two kinds of IPC. The mechanisms with a real
security story---seL4's capabilities, CHERI-backed isolation---demand a
new kernel or new silicon, neither of which can be rolled out to an
existing fleet. The mechanisms that are actually fast on stock
Linux---Aeron, io\_uring, eRPC---will happily deliver a message from
anyone to anyone.

\sysnamefull closes that gap. It is, to our knowledge, the first
userspace runtime to put capability-token access control on a
zero-copy memfd IPC fastpath on stock Linux: no kernel patch, no
special hardware, no hypervisor. Each endpoint of every channel holds
an unforgeable 128-bit token, and every message---once at the
producer, again at the consumer---must pass validation before it
moves.

The check is affordable because of a \emph{slot-indexed validate
path}, which collapses what would naturally be an $O(n)$ pool scan
into a single shift, one cache load, and one constant-time compare,
while keeping a $2^{96}$ forgery space. The gate costs
\gateOverheadTail\,ns per message at p99---\gateOverheadPercent\% of
an ungated round trip that is itself only \rawAstraRTT\,ns---and the
resulting mediated path still beats every transport we measured
against, Aeron included. On top of that it provides
Cornucopia\,Reloaded-class cascading revocation in software, with
$O(1)$ time-to-effect regardless of pool occupancy. Revocation lands
within \revokeTail\,$\mu$s of the \texttt{revoke()} call, end to end
including drain; enabling the gate costs \mpscGateCostPercent\% of
single-producer MPSC throughput; and the implementation is checked
against a hand-written
property-based oracle and a CBMC bounded-model proof of permission
monotonicity.

Our evaluation runs pipe(2), socketpair, io\_uring, and Aeron
head-to-head, all in the same TSC domain on one machine. Source,
baseline harnesses, and a one-command artefact-evaluation sweep
accompany the paper.
